\chapter{Resultados y Análisis}
\label{cap:resultados}

En este capítulo se presentan y evalúan los resultados obtenidos tras la ejecución del simulador BCH desarrollado. El análisis se centra en validar la eficiencia de las optimizaciones implementadas (como el \textit{bit-packing}) y la precisión matemática del sistema en la corrección de errores sobre canales con ruido.

\section{Entorno de Experimentación}
Para garantizar la reproducibilidad de las pruebas, se describen a continuación las características del hardware y el software utilizados durante las simulaciones de Monte Carlo.

\subsection{Configuración del Host (Hardware Base)}
Aunque la ejecución se realiza en un entorno controlado, las capacidades del hardware físico subyacente definen el rendimiento máximo alcanzable:
\begin{itemize}
    \item \textbf{Procesador Host:} AMD Ryzen AI 9 HX 370 (12 núcleos, 24 hilos), que proporciona una alta capacidad de cómputo para las iteraciones intensivas.
    \item \textbf{Memoria Física:} 32 GB LPDDR5, permitiendo el manejo fluido de grandes volúmenes de datos en memoria.
    \item \textbf{GPU:} NVIDIA RTX 5070, que aunque no interviene directamente en el motor de ejecución serie, sitúa el desarrollo en una arquitectura de alto rendimiento.
\end{itemize}

\subsection{Entorno de Virtualización}
Debido a necesidades de aislamiento del entorno de desarrollo, las pruebas se han ejecutado en un sistema virtualizado, lo cual introduce un \textit{overhead} que debe considerarse en el análisis de tiempos:
\begin{itemize}
    \item \textbf{Hipervisor:} Oracle VM VirtualBox 7.x.
    \item \textbf{Sistema Operativo Invitado (Guest):} [Insertar tu SO, ej: Ubuntu 22.04 LTS], configurado con [X] vCPUs y [Y] GB de memoria RAM.
    \item \textbf{Nota sobre el rendimiento:} Dado que el análisis de escalabilidad es comparativo entre diferentes valores de $m$, las conclusiones sobre la eficiencia relativa de los algoritmos siguen siendo plenamente válidas a pesar de la capa de virtualización.
\end{itemize}

\subsection{Configuración de Software}
\begin{itemize}
    \item \textbf{Compilador:} GCC con el estándar C++17, utilizando el flag de optimización \texttt{-O3} para maximizar el rendimiento de las operaciones a nivel de bit.
    \item \textbf{Post-procesado:} Scripts en Python 3 utilizando las librerías \textit{Pandas} para la gestión de datos CSV y \textit{Matplotlib} para la generación de curvas de error y tiempos.
\end{itemize}

\section{Metodología de Validación}
Para obtener resultados estadísticamente significativos, cada punto de la relación señal-ruido ($E_b/N_0$) se ha simulado bajo un criterio de parada riguroso: la simulación continúa hasta contabilizar un mínimo de 100 tramas erróneas (FER) o hasta alcanzar un límite de $10^6$ tramas transmitidas. Este método asegura que las curvas de BER converjan y sean representativas del comportamiento teórico del código.

\section{Análisis de Escalabilidad Temporal}
En esta sección se evalúa el impacto del tamaño del campo $GF(2^m)$ en los tiempos de proceso.

\begin{figure}[h]
    \centering
    \includegraphics[width=0.8\textwidth]{results/plot/grafica_escalabilidad_m.png}
    \caption{Tiempos promedio de codificación y decodificación en función de $m$.}
    \label{fig:grafica_tiempos}
\end{figure}

\subsection{Impacto del Bit-packing en la Codificación}
Como se observa en la Figura \ref{fig:grafica_tiempos}, el tiempo de codificación se mantiene notablemente bajo y con un crecimiento moderado a pesar del aumento exponencial de $n$. Esto confirma la eficacia de la estructura \texttt{PackedBit} y el método \texttt{xorShifted}, que permiten procesar 64 bits por cada operación de ciclo de CPU virtualizado.

\subsection{Complejidad de la Decodificación}
El tiempo de decodificación muestra un crecimiento más pronunciado, lo cual es coherente con la complejidad algorítmica de Berlekamp-Massey y la búsqueda de Chien. No obstante, el uso de \texttt{uint16_t} para las tablas de búsqueda asegura que la latencia de acceso a memoria se mantenga mínima al favorecer la localidad de datos en la caché.

\section{Evaluación del Rendimiento (BER y FER)}
El análisis de fiabilidad constituye el núcleo de la evaluación del sistema.

\begin{figure}[h]
    \centering
    \includegraphics[width=0.8\textwidth]{results/plot/grafica_ber_m15.png}
    \caption{Curvas de BER y FER para un código BCH ($m=15$).}
    \label{fig:grafica_ber}
\end{figure}

\subsection{Ganancia de Codificación}
La Figura \ref{fig:grafica_ber} compara la tasa de error de bit (BER) obtenida con el simulador frente a un sistema sin codificar (\textit{Uncoded}). La distancia horizontal entre ambas curvas representa la ganancia de codificación, demostrando la capacidad del simulador para corregir hasta $t$ errores por bloque.

\section{Discusión de Resultados}
Los resultados obtenidos confirman que la implementación es matemáticamente robusta. La elección de una arquitectura de 16 bits no solo ha optimizado el uso de memoria, sino que ha permitido alcanzar longitudes de bloque de hasta $n=32.767$ con tiempos de respuesta aceptables en un entorno virtualizado.